% !TeX program = lualatex
% !TeX root = predavanje04.tex
% ============================================================
% Zajednički preamble za sva predavanja RWA kolegija
% Uključiti s: % ============================================================
% Zajednički preamble za sva predavanja RWA kolegija
% Uključiti s: % ============================================================
% Zajednički preamble za sva predavanja RWA kolegija
% Uključiti s: \input{../common/preamble.tex}
% ============================================================

\documentclass[aspectratio=149,11pt]{beamer}
\usetheme{CambridgeUS}

\usepackage{fontspec}              % radi samo sa XeLaTeX/LuaLaTeX
\usepackage[croatian]{babel}
\usepackage{csquotes}              % \enquote{} → hrvatski navodnici „..."
\usepackage{booktabs}
\usepackage{array}
\usepackage{listings}
\usepackage{xcolor}
\usepackage{colortbl}
\usepackage{hyperref}
\usepackage{tikz}
\usetikzlibrary{positioning,arrows.meta}

% ---------- Tipografija: IBM Plex ----------
\setsansfont{IBMPlexSans}[
  Extension      = .otf,
  UprightFont    = *-Regular,
  ItalicFont     = *-Italic,
  BoldFont       = *-SemiBold,
  BoldItalicFont = *-SemiBoldItalic
]
\setmainfont{IBMPlexSans}[
  Extension      = .otf,
  UprightFont    = *-Regular,
  ItalicFont     = *-Italic,
  BoldFont       = *-SemiBold,
  BoldItalicFont = *-SemiBoldItalic
]
\setmonofont{IBMPlexMono}[
  Extension      = .otf,
  UprightFont    = *-Regular,
  ItalicFont     = *-Italic,
  BoldFont       = *-SemiBold,
  BoldItalicFont = *-SemiBoldItalic,
  Scale          = 0.92
]

% ---------- Paleta ----------
\definecolor{ink}{HTML}{0C4685}      % primarna tamnoplava
\definecolor{ink2}{HTML}{6E7BAA}     % sekundarna plavo-siva
\definecolor{tint}{HTML}{EAEEF6}     % svjetli plavi tint
\definecolor{textclr}{HTML}{1A222E}  % glavni tekst
\definecolor{mutedclr}{HTML}{6B7484} % podnožje, sekundarni tekst
\definecolor{bgclr}{HTML}{F6F7F9}    % pozadina slajda
\definecolor{lineclr}{HTML}{D9DEE7}  % tanke linije tablica
\definecolor{divrule}{HTML}{9DB4D6}  % akcent na tamnim sekcijskim slajdovima
\definecolor{divtext}{HTML}{E6ECF6}  % tekst na tamnim sekcijskim slajdovima

\setbeamercolor{normal text}{fg=textclr,bg=bgclr}
\setbeamercolor{structure}{fg=ink}
\setbeamercolor{alerted text}{fg=ink}

\setbeamercolor{title}{fg=ink,bg=tint}
\setbeamercolor{subtitle}{fg=mutedclr}
\setbeamercolor{author}{fg=textclr}
\setbeamercolor{date}{fg=mutedclr}
\setbeamercolor{frametitle}{fg=ink,bg=tint}

\setbeamercolor{section in head/foot}{fg=white,bg=ink}
\setbeamercolor{subsection in head/foot}{fg=ink,bg=tint}
\setbeamercolor{author in head/foot}{fg=white,bg=ink}
\setbeamercolor{title in head/foot}{fg=textclr,bg=tint}
\setbeamercolor{date in head/foot}{fg=mutedclr,bg=tint}

\setbeamercolor{block title}{fg=white,bg=ink}
\setbeamercolor{block body}{fg=textclr,bg=tint}
\setbeamercolor{block title example}{fg=white,bg=ink2}
\setbeamercolor{block body example}{fg=textclr,bg=tint}
\setbeamercolor{itemize item}{fg=ink2}
\setbeamercolor{itemize subitem}{fg=ink2}
\setbeamercolor{enumerate item}{fg=ink2}

\setbeamertemplate{itemize items}[square]
\setbeamertemplate{navigation symbols}{}
\arrayrulecolor{lineclr}

% ---------- Kolegij, autori i logo ----------
\newcommand{\kolegij}{Programiranje za web}
\author{doc.\,dr.\,sc. Tomislav Petković \and mag.\,ing. Josip Torić}
\newcommand{\znak}[2]{\raisebox{-0.5\height}{\includegraphics[height=#1]{#2}}}
\titlegraphic{\znak{2.1cm}{../common/SIT_logo_tekst1.pdf}\hspace{1.8em}\znak{1.8cm}{../common/UNIZD_logo.pdf}}

% ---------- Zaglavlje ----------
% Jedna puna traka: lijevo sekcija, desno podsekcija.
\setbeamertemplate{headline}{%
  \leavevmode%
  \hbox{%
    \begin{beamercolorbox}[wd=\paperwidth,ht=2.25ex,dp=1ex,leftskip=.3cm,rightskip=.3cm]{section in head/foot}%
      \usebeamerfont{section in head/foot}\insertsectionhead\hfill\insertsubsectionhead%
    \end{beamercolorbox}%
  }%
}

% ---------- Podnožje ----------
% Lijevo: naziv trenutne prezentacije (kratki oblik naslova iz \title[...]{...})
% Sredina: naziv kolegija.  Desno: datum i redni broj slajda.
\setbeamertemplate{footline}{%
  \leavevmode%
  \hbox{%
    \begin{beamercolorbox}[wd=.34\paperwidth,ht=2.25ex,dp=1ex,center]{author in head/foot}%
      \usebeamerfont{author in head/foot}\insertshorttitle%
    \end{beamercolorbox}%
    \begin{beamercolorbox}[wd=.33\paperwidth,ht=2.25ex,dp=1ex,center]{title in head/foot}%
      \usebeamerfont{title in head/foot}\kolegij%
    \end{beamercolorbox}%
    \begin{beamercolorbox}[wd=.33\paperwidth,ht=2.25ex,dp=1ex,center]{date in head/foot}%
      \usebeamerfont{date in head/foot}\insertshortdate{}\hspace*{2em}%
      \insertframenumber{}\,/\,\inserttotalframenumber\hspace*{1ex}%
    \end{beamercolorbox}%
  }%
  \vskip0pt%
}

% ---------- Sekcijski slajd ----------
% \sekcija{Naslov sekcije}{jedna linija koja nabraja teme sekcije}
\newcommand{\sekcija}[2]{%
  \section{#1}%
  {%
    \setbeamercolor{background canvas}{bg=ink}%
    \setbeamertemplate{headline}{}%
    \begin{frame}
      \vfill
      \hspace*{0.03\textwidth}{\color{divrule}\rule{2.6cm}{1.1pt}}\par
      \vskip 0.7em
      \hspace*{0.03\textwidth}{\Large\bfseries\color{white}#1}\par
      \vskip 0.6em
      \hspace*{0.03\textwidth}\begin{minipage}{0.9\textwidth}\color{divtext}\small #2\end{minipage}\par
      \vfill
    \end{frame}%
  }%
}

% ---------- Listings ----------
\lstdefinelanguage{bash}{
  morekeywords={sudo,apt,brew,pip,pipx,uv,winget,python,git,flask,source,deactivate,code,gunicorn,docker,compose,heroku,render,railway,node,npm,pnpm,yarn,uvicorn,alembic,pytest,ruff,black,precommit},
  sensitive=false,
  morecomment=[l]{\#}
}
\lstdefinelanguage{python}{
  morekeywords={from,import,def,class,return,if,elif,else,for,in,while,try,except,raise,with,as,not,and,or,is,lambda,yield,True,False,None,await,async},
  sensitive=true,
  morecomment=[l]{\#},
  morestring=[b]',
  morestring=[b]"
}
\lstdefinelanguage{sql}{
  morekeywords={SELECT,FROM,WHERE,INSERT,INTO,VALUES,UPDATE,SET,DELETE,CREATE,TABLE,PRIMARY,KEY,FOREIGN,REFERENCES,UNIQUE,NOT,NULL,INDEX,ON,ALTER,DROP,CONSTRAINT},
  sensitive=false,
  morecomment=[l]{--}
}
\lstdefinelanguage{json}{
  morecomment=[l]{//},
  morestring=[b]",
  showstringspaces=false
}
\lstdefinelanguage{typescript}{
  morekeywords={import,export,from,const,let,var,function,return,if,else,for,of,in,while,switch,case,break,async,await,try,catch,finally,throw,new,class,this,extends,implements,interface,type,enum,readonly,public,private,protected,static,get,set,true,false,null,undefined,void,string,number,boolean,any,unknown,never,typeof,keyof,as},
  sensitive=true,
  morecomment=[l]{//},
  morecomment=[s]{/*}{*/},
  morestring=[b]',
  morestring=[b]",
  morestring=[b]`
}
\lstalias{javascript}{typescript}
\lstset{
  backgroundcolor=\color{tint},
  basicstyle=\ttfamily\footnotesize,
  keywordstyle=\color{ink}\bfseries,
  commentstyle=\color{mutedclr},
  frame=single,
  rulecolor=\color{lineclr},
  frameround=tttt,
  breaklines=true,
  tabsize=2,
  showstringspaces=false,
  upquote=true,
  columns=fullflexible,
  escapeinside={<@}{@>}
}
% ---------- Dijakritici u isječcima koda ----------
% listings čita bajt po bajt, pa višebajtni znak premjesti: "početni" ispadne
% "čpoetni". literate to ne popravlja. Rješenje je izlaz u LaTeX: sve između
% <@ i @> ne prolazi kroz listings, nego se slaže kao običan tekst, a unutar
% okvira je font ionako ttfamily, pa se ništa ne razlikuje.
%
%   async def get(...):
%       club = await repo.get_by_id(db, club_id)   <@\kom{dohvat kluba}@>
%       raise AppError("not_found", <@"Klub nije pronađen"@>, 404)
%
% Graničnici <@ i @> odabrani su jer se ne pojavljuju ni u jednom jeziku
% koji se na predavanjima citira; (* bi se sudaralo s def f(*args).
\newcommand{\kom}[1]{{\color{mutedclr}\# #1}}
% Isto za TypeScript, Vue i CSS, kojima komentar počinje s //:
%   const lifters = useLiftersStore()   <@\komts{dijeljeni popis}@>
\newcommand{\komts}[1]{{\color{mutedclr}// #1}}

% ---------- Činjenica naspram konvencije ----------
% Sve što nije propisano standardom ili jezikom, nego je dogovor struke,
% ide u ovaj okvir, da student zna koje je vrste tvrdnja pred njim.
% \begin{konvencija} ... \end{konvencija}
% \begin{konvencija}[Konvencija: nazivi putanja] ... \end{konvencija}
\newenvironment{konvencija}[1][Konvencija]
  {\begin{exampleblock}{#1}}
  {\end{exampleblock}}

% ---------- Helper macros ----------
% Stupac tablice zadane širine, poravnat ulijevo (bez rastezanja i dijeljenja riječi).
\newcolumntype{L}[1]{>{\raggedright\arraybackslash}p{#1}}
\newcommand{\codefile}[1]{\texttt{#1}}
\newcommand{\kbd}[1]{\texttt{#1}}

% ---------- Code frame environment ----------
% Automatski dodaje fragile i shrink=5, koristi se umjesto \begin{frame}
% Primjer: \begin{codeframe}{Naslov} ... \end{codeframe}
% Primjer s opcijama: \begin{codeframe}[shrink=15]{Naslov} ... \end{codeframe}
\newenvironment{codeframe}[2][]
  {\begin{frame}[fragile,shrink=5,#1]{#2}}
  {\end{frame}}

% ============================================================

\documentclass[aspectratio=149,11pt]{beamer}
\usetheme{CambridgeUS}

\usepackage{fontspec}              % radi samo sa XeLaTeX/LuaLaTeX
\usepackage[croatian]{babel}
\usepackage{csquotes}              % \enquote{} → hrvatski navodnici „..."
\usepackage{booktabs}
\usepackage{array}
\usepackage{listings}
\usepackage{xcolor}
\usepackage{colortbl}
\usepackage{hyperref}
\usepackage{tikz}
\usetikzlibrary{positioning,arrows.meta}

% ---------- Tipografija: IBM Plex ----------
\setsansfont{IBMPlexSans}[
  Extension      = .otf,
  UprightFont    = *-Regular,
  ItalicFont     = *-Italic,
  BoldFont       = *-SemiBold,
  BoldItalicFont = *-SemiBoldItalic
]
\setmainfont{IBMPlexSans}[
  Extension      = .otf,
  UprightFont    = *-Regular,
  ItalicFont     = *-Italic,
  BoldFont       = *-SemiBold,
  BoldItalicFont = *-SemiBoldItalic
]
\setmonofont{IBMPlexMono}[
  Extension      = .otf,
  UprightFont    = *-Regular,
  ItalicFont     = *-Italic,
  BoldFont       = *-SemiBold,
  BoldItalicFont = *-SemiBoldItalic,
  Scale          = 0.92
]

% ---------- Paleta ----------
\definecolor{ink}{HTML}{0C4685}      % primarna tamnoplava
\definecolor{ink2}{HTML}{6E7BAA}     % sekundarna plavo-siva
\definecolor{tint}{HTML}{EAEEF6}     % svjetli plavi tint
\definecolor{textclr}{HTML}{1A222E}  % glavni tekst
\definecolor{mutedclr}{HTML}{6B7484} % podnožje, sekundarni tekst
\definecolor{bgclr}{HTML}{F6F7F9}    % pozadina slajda
\definecolor{lineclr}{HTML}{D9DEE7}  % tanke linije tablica
\definecolor{divrule}{HTML}{9DB4D6}  % akcent na tamnim sekcijskim slajdovima
\definecolor{divtext}{HTML}{E6ECF6}  % tekst na tamnim sekcijskim slajdovima

\setbeamercolor{normal text}{fg=textclr,bg=bgclr}
\setbeamercolor{structure}{fg=ink}
\setbeamercolor{alerted text}{fg=ink}

\setbeamercolor{title}{fg=ink,bg=tint}
\setbeamercolor{subtitle}{fg=mutedclr}
\setbeamercolor{author}{fg=textclr}
\setbeamercolor{date}{fg=mutedclr}
\setbeamercolor{frametitle}{fg=ink,bg=tint}

\setbeamercolor{section in head/foot}{fg=white,bg=ink}
\setbeamercolor{subsection in head/foot}{fg=ink,bg=tint}
\setbeamercolor{author in head/foot}{fg=white,bg=ink}
\setbeamercolor{title in head/foot}{fg=textclr,bg=tint}
\setbeamercolor{date in head/foot}{fg=mutedclr,bg=tint}

\setbeamercolor{block title}{fg=white,bg=ink}
\setbeamercolor{block body}{fg=textclr,bg=tint}
\setbeamercolor{block title example}{fg=white,bg=ink2}
\setbeamercolor{block body example}{fg=textclr,bg=tint}
\setbeamercolor{itemize item}{fg=ink2}
\setbeamercolor{itemize subitem}{fg=ink2}
\setbeamercolor{enumerate item}{fg=ink2}

\setbeamertemplate{itemize items}[square]
\setbeamertemplate{navigation symbols}{}
\arrayrulecolor{lineclr}

% ---------- Kolegij, autori i logo ----------
\newcommand{\kolegij}{Programiranje za web}
\author{doc.\,dr.\,sc. Tomislav Petković \and mag.\,ing. Josip Torić}
\newcommand{\znak}[2]{\raisebox{-0.5\height}{\includegraphics[height=#1]{#2}}}
\titlegraphic{\znak{2.1cm}{../common/SIT_logo_tekst1.pdf}\hspace{1.8em}\znak{1.8cm}{../common/UNIZD_logo.pdf}}

% ---------- Zaglavlje ----------
% Jedna puna traka: lijevo sekcija, desno podsekcija.
\setbeamertemplate{headline}{%
  \leavevmode%
  \hbox{%
    \begin{beamercolorbox}[wd=\paperwidth,ht=2.25ex,dp=1ex,leftskip=.3cm,rightskip=.3cm]{section in head/foot}%
      \usebeamerfont{section in head/foot}\insertsectionhead\hfill\insertsubsectionhead%
    \end{beamercolorbox}%
  }%
}

% ---------- Podnožje ----------
% Lijevo: naziv trenutne prezentacije (kratki oblik naslova iz \title[...]{...})
% Sredina: naziv kolegija.  Desno: datum i redni broj slajda.
\setbeamertemplate{footline}{%
  \leavevmode%
  \hbox{%
    \begin{beamercolorbox}[wd=.34\paperwidth,ht=2.25ex,dp=1ex,center]{author in head/foot}%
      \usebeamerfont{author in head/foot}\insertshorttitle%
    \end{beamercolorbox}%
    \begin{beamercolorbox}[wd=.33\paperwidth,ht=2.25ex,dp=1ex,center]{title in head/foot}%
      \usebeamerfont{title in head/foot}\kolegij%
    \end{beamercolorbox}%
    \begin{beamercolorbox}[wd=.33\paperwidth,ht=2.25ex,dp=1ex,center]{date in head/foot}%
      \usebeamerfont{date in head/foot}\insertshortdate{}\hspace*{2em}%
      \insertframenumber{}\,/\,\inserttotalframenumber\hspace*{1ex}%
    \end{beamercolorbox}%
  }%
  \vskip0pt%
}

% ---------- Sekcijski slajd ----------
% \sekcija{Naslov sekcije}{jedna linija koja nabraja teme sekcije}
\newcommand{\sekcija}[2]{%
  \section{#1}%
  {%
    \setbeamercolor{background canvas}{bg=ink}%
    \setbeamertemplate{headline}{}%
    \begin{frame}
      \vfill
      \hspace*{0.03\textwidth}{\color{divrule}\rule{2.6cm}{1.1pt}}\par
      \vskip 0.7em
      \hspace*{0.03\textwidth}{\Large\bfseries\color{white}#1}\par
      \vskip 0.6em
      \hspace*{0.03\textwidth}\begin{minipage}{0.9\textwidth}\color{divtext}\small #2\end{minipage}\par
      \vfill
    \end{frame}%
  }%
}

% ---------- Listings ----------
\lstdefinelanguage{bash}{
  morekeywords={sudo,apt,brew,pip,pipx,uv,winget,python,git,flask,source,deactivate,code,gunicorn,docker,compose,heroku,render,railway,node,npm,pnpm,yarn,uvicorn,alembic,pytest,ruff,black,precommit},
  sensitive=false,
  morecomment=[l]{\#}
}
\lstdefinelanguage{python}{
  morekeywords={from,import,def,class,return,if,elif,else,for,in,while,try,except,raise,with,as,not,and,or,is,lambda,yield,True,False,None,await,async},
  sensitive=true,
  morecomment=[l]{\#},
  morestring=[b]',
  morestring=[b]"
}
\lstdefinelanguage{sql}{
  morekeywords={SELECT,FROM,WHERE,INSERT,INTO,VALUES,UPDATE,SET,DELETE,CREATE,TABLE,PRIMARY,KEY,FOREIGN,REFERENCES,UNIQUE,NOT,NULL,INDEX,ON,ALTER,DROP,CONSTRAINT},
  sensitive=false,
  morecomment=[l]{--}
}
\lstdefinelanguage{json}{
  morecomment=[l]{//},
  morestring=[b]",
  showstringspaces=false
}
\lstdefinelanguage{typescript}{
  morekeywords={import,export,from,const,let,var,function,return,if,else,for,of,in,while,switch,case,break,async,await,try,catch,finally,throw,new,class,this,extends,implements,interface,type,enum,readonly,public,private,protected,static,get,set,true,false,null,undefined,void,string,number,boolean,any,unknown,never,typeof,keyof,as},
  sensitive=true,
  morecomment=[l]{//},
  morecomment=[s]{/*}{*/},
  morestring=[b]',
  morestring=[b]",
  morestring=[b]`
}
\lstalias{javascript}{typescript}
\lstset{
  backgroundcolor=\color{tint},
  basicstyle=\ttfamily\footnotesize,
  keywordstyle=\color{ink}\bfseries,
  commentstyle=\color{mutedclr},
  frame=single,
  rulecolor=\color{lineclr},
  frameround=tttt,
  breaklines=true,
  tabsize=2,
  showstringspaces=false,
  upquote=true,
  columns=fullflexible,
  escapeinside={<@}{@>}
}
% ---------- Dijakritici u isječcima koda ----------
% listings čita bajt po bajt, pa višebajtni znak premjesti: "početni" ispadne
% "čpoetni". literate to ne popravlja. Rješenje je izlaz u LaTeX: sve između
% <@ i @> ne prolazi kroz listings, nego se slaže kao običan tekst, a unutar
% okvira je font ionako ttfamily, pa se ništa ne razlikuje.
%
%   async def get(...):
%       club = await repo.get_by_id(db, club_id)   <@\kom{dohvat kluba}@>
%       raise AppError("not_found", <@"Klub nije pronađen"@>, 404)
%
% Graničnici <@ i @> odabrani su jer se ne pojavljuju ni u jednom jeziku
% koji se na predavanjima citira; (* bi se sudaralo s def f(*args).
\newcommand{\kom}[1]{{\color{mutedclr}\# #1}}
% Isto za TypeScript, Vue i CSS, kojima komentar počinje s //:
%   const lifters = useLiftersStore()   <@\komts{dijeljeni popis}@>
\newcommand{\komts}[1]{{\color{mutedclr}// #1}}

% ---------- Činjenica naspram konvencije ----------
% Sve što nije propisano standardom ili jezikom, nego je dogovor struke,
% ide u ovaj okvir, da student zna koje je vrste tvrdnja pred njim.
% \begin{konvencija} ... \end{konvencija}
% \begin{konvencija}[Konvencija: nazivi putanja] ... \end{konvencija}
\newenvironment{konvencija}[1][Konvencija]
  {\begin{exampleblock}{#1}}
  {\end{exampleblock}}

% ---------- Helper macros ----------
% Stupac tablice zadane širine, poravnat ulijevo (bez rastezanja i dijeljenja riječi).
\newcolumntype{L}[1]{>{\raggedright\arraybackslash}p{#1}}
\newcommand{\codefile}[1]{\texttt{#1}}
\newcommand{\kbd}[1]{\texttt{#1}}

% ---------- Code frame environment ----------
% Automatski dodaje fragile i shrink=5, koristi se umjesto \begin{frame}
% Primjer: \begin{codeframe}{Naslov} ... \end{codeframe}
% Primjer s opcijama: \begin{codeframe}[shrink=15]{Naslov} ... \end{codeframe}
\newenvironment{codeframe}[2][]
  {\begin{frame}[fragile,shrink=5,#1]{#2}}
  {\end{frame}}

% ============================================================

\documentclass[aspectratio=149,11pt]{beamer}
\usetheme{CambridgeUS}

\usepackage{fontspec}              % radi samo sa XeLaTeX/LuaLaTeX
\usepackage[croatian]{babel}
\usepackage{csquotes}              % \enquote{} → hrvatski navodnici „..."
\usepackage{booktabs}
\usepackage{array}
\usepackage{listings}
\usepackage{xcolor}
\usepackage{colortbl}
\usepackage{hyperref}
\usepackage{tikz}
\usetikzlibrary{positioning,arrows.meta}

% ---------- Tipografija: IBM Plex ----------
\setsansfont{IBMPlexSans}[
  Extension      = .otf,
  UprightFont    = *-Regular,
  ItalicFont     = *-Italic,
  BoldFont       = *-SemiBold,
  BoldItalicFont = *-SemiBoldItalic
]
\setmainfont{IBMPlexSans}[
  Extension      = .otf,
  UprightFont    = *-Regular,
  ItalicFont     = *-Italic,
  BoldFont       = *-SemiBold,
  BoldItalicFont = *-SemiBoldItalic
]
\setmonofont{IBMPlexMono}[
  Extension      = .otf,
  UprightFont    = *-Regular,
  ItalicFont     = *-Italic,
  BoldFont       = *-SemiBold,
  BoldItalicFont = *-SemiBoldItalic,
  Scale          = 0.92
]

% ---------- Paleta ----------
\definecolor{ink}{HTML}{0C4685}      % primarna tamnoplava
\definecolor{ink2}{HTML}{6E7BAA}     % sekundarna plavo-siva
\definecolor{tint}{HTML}{EAEEF6}     % svjetli plavi tint
\definecolor{textclr}{HTML}{1A222E}  % glavni tekst
\definecolor{mutedclr}{HTML}{6B7484} % podnožje, sekundarni tekst
\definecolor{bgclr}{HTML}{F6F7F9}    % pozadina slajda
\definecolor{lineclr}{HTML}{D9DEE7}  % tanke linije tablica
\definecolor{divrule}{HTML}{9DB4D6}  % akcent na tamnim sekcijskim slajdovima
\definecolor{divtext}{HTML}{E6ECF6}  % tekst na tamnim sekcijskim slajdovima

\setbeamercolor{normal text}{fg=textclr,bg=bgclr}
\setbeamercolor{structure}{fg=ink}
\setbeamercolor{alerted text}{fg=ink}

\setbeamercolor{title}{fg=ink,bg=tint}
\setbeamercolor{subtitle}{fg=mutedclr}
\setbeamercolor{author}{fg=textclr}
\setbeamercolor{date}{fg=mutedclr}
\setbeamercolor{frametitle}{fg=ink,bg=tint}

\setbeamercolor{section in head/foot}{fg=white,bg=ink}
\setbeamercolor{subsection in head/foot}{fg=ink,bg=tint}
\setbeamercolor{author in head/foot}{fg=white,bg=ink}
\setbeamercolor{title in head/foot}{fg=textclr,bg=tint}
\setbeamercolor{date in head/foot}{fg=mutedclr,bg=tint}

\setbeamercolor{block title}{fg=white,bg=ink}
\setbeamercolor{block body}{fg=textclr,bg=tint}
\setbeamercolor{block title example}{fg=white,bg=ink2}
\setbeamercolor{block body example}{fg=textclr,bg=tint}
\setbeamercolor{itemize item}{fg=ink2}
\setbeamercolor{itemize subitem}{fg=ink2}
\setbeamercolor{enumerate item}{fg=ink2}

\setbeamertemplate{itemize items}[square]
\setbeamertemplate{navigation symbols}{}
\arrayrulecolor{lineclr}

% ---------- Kolegij, autori i logo ----------
\newcommand{\kolegij}{Programiranje za web}
\author{doc.\,dr.\,sc. Tomislav Petković \and mag.\,ing. Josip Torić}
\newcommand{\znak}[2]{\raisebox{-0.5\height}{\includegraphics[height=#1]{#2}}}
\titlegraphic{\znak{2.1cm}{../common/SIT_logo_tekst1.pdf}\hspace{1.8em}\znak{1.8cm}{../common/UNIZD_logo.pdf}}

% ---------- Zaglavlje ----------
% Jedna puna traka: lijevo sekcija, desno podsekcija.
\setbeamertemplate{headline}{%
  \leavevmode%
  \hbox{%
    \begin{beamercolorbox}[wd=\paperwidth,ht=2.25ex,dp=1ex,leftskip=.3cm,rightskip=.3cm]{section in head/foot}%
      \usebeamerfont{section in head/foot}\insertsectionhead\hfill\insertsubsectionhead%
    \end{beamercolorbox}%
  }%
}

% ---------- Podnožje ----------
% Lijevo: naziv trenutne prezentacije (kratki oblik naslova iz \title[...]{...})
% Sredina: naziv kolegija.  Desno: datum i redni broj slajda.
\setbeamertemplate{footline}{%
  \leavevmode%
  \hbox{%
    \begin{beamercolorbox}[wd=.34\paperwidth,ht=2.25ex,dp=1ex,center]{author in head/foot}%
      \usebeamerfont{author in head/foot}\insertshorttitle%
    \end{beamercolorbox}%
    \begin{beamercolorbox}[wd=.33\paperwidth,ht=2.25ex,dp=1ex,center]{title in head/foot}%
      \usebeamerfont{title in head/foot}\kolegij%
    \end{beamercolorbox}%
    \begin{beamercolorbox}[wd=.33\paperwidth,ht=2.25ex,dp=1ex,center]{date in head/foot}%
      \usebeamerfont{date in head/foot}\insertshortdate{}\hspace*{2em}%
      \insertframenumber{}\,/\,\inserttotalframenumber\hspace*{1ex}%
    \end{beamercolorbox}%
  }%
  \vskip0pt%
}

% ---------- Sekcijski slajd ----------
% \sekcija{Naslov sekcije}{jedna linija koja nabraja teme sekcije}
\newcommand{\sekcija}[2]{%
  \section{#1}%
  {%
    \setbeamercolor{background canvas}{bg=ink}%
    \setbeamertemplate{headline}{}%
    \begin{frame}
      \vfill
      \hspace*{0.03\textwidth}{\color{divrule}\rule{2.6cm}{1.1pt}}\par
      \vskip 0.7em
      \hspace*{0.03\textwidth}{\Large\bfseries\color{white}#1}\par
      \vskip 0.6em
      \hspace*{0.03\textwidth}\begin{minipage}{0.9\textwidth}\color{divtext}\small #2\end{minipage}\par
      \vfill
    \end{frame}%
  }%
}

% ---------- Listings ----------
\lstdefinelanguage{bash}{
  morekeywords={sudo,apt,brew,pip,pipx,uv,winget,python,git,flask,source,deactivate,code,gunicorn,docker,compose,heroku,render,railway,node,npm,pnpm,yarn,uvicorn,alembic,pytest,ruff,black,precommit},
  sensitive=false,
  morecomment=[l]{\#}
}
\lstdefinelanguage{python}{
  morekeywords={from,import,def,class,return,if,elif,else,for,in,while,try,except,raise,with,as,not,and,or,is,lambda,yield,True,False,None,await,async},
  sensitive=true,
  morecomment=[l]{\#},
  morestring=[b]',
  morestring=[b]"
}
\lstdefinelanguage{sql}{
  morekeywords={SELECT,FROM,WHERE,INSERT,INTO,VALUES,UPDATE,SET,DELETE,CREATE,TABLE,PRIMARY,KEY,FOREIGN,REFERENCES,UNIQUE,NOT,NULL,INDEX,ON,ALTER,DROP,CONSTRAINT},
  sensitive=false,
  morecomment=[l]{--}
}
\lstdefinelanguage{json}{
  morecomment=[l]{//},
  morestring=[b]",
  showstringspaces=false
}
\lstdefinelanguage{typescript}{
  morekeywords={import,export,from,const,let,var,function,return,if,else,for,of,in,while,switch,case,break,async,await,try,catch,finally,throw,new,class,this,extends,implements,interface,type,enum,readonly,public,private,protected,static,get,set,true,false,null,undefined,void,string,number,boolean,any,unknown,never,typeof,keyof,as},
  sensitive=true,
  morecomment=[l]{//},
  morecomment=[s]{/*}{*/},
  morestring=[b]',
  morestring=[b]",
  morestring=[b]`
}
\lstalias{javascript}{typescript}
\lstset{
  backgroundcolor=\color{tint},
  basicstyle=\ttfamily\footnotesize,
  keywordstyle=\color{ink}\bfseries,
  commentstyle=\color{mutedclr},
  frame=single,
  rulecolor=\color{lineclr},
  frameround=tttt,
  breaklines=true,
  tabsize=2,
  showstringspaces=false,
  upquote=true,
  columns=fullflexible,
  escapeinside={<@}{@>}
}
% ---------- Dijakritici u isječcima koda ----------
% listings čita bajt po bajt, pa višebajtni znak premjesti: "početni" ispadne
% "čpoetni". literate to ne popravlja. Rješenje je izlaz u LaTeX: sve između
% <@ i @> ne prolazi kroz listings, nego se slaže kao običan tekst, a unutar
% okvira je font ionako ttfamily, pa se ništa ne razlikuje.
%
%   async def get(...):
%       club = await repo.get_by_id(db, club_id)   <@\kom{dohvat kluba}@>
%       raise AppError("not_found", <@"Klub nije pronađen"@>, 404)
%
% Graničnici <@ i @> odabrani su jer se ne pojavljuju ni u jednom jeziku
% koji se na predavanjima citira; (* bi se sudaralo s def f(*args).
\newcommand{\kom}[1]{{\color{mutedclr}\# #1}}
% Isto za TypeScript, Vue i CSS, kojima komentar počinje s //:
%   const lifters = useLiftersStore()   <@\komts{dijeljeni popis}@>
\newcommand{\komts}[1]{{\color{mutedclr}// #1}}

% ---------- Činjenica naspram konvencije ----------
% Sve što nije propisano standardom ili jezikom, nego je dogovor struke,
% ide u ovaj okvir, da student zna koje je vrste tvrdnja pred njim.
% \begin{konvencija} ... \end{konvencija}
% \begin{konvencija}[Konvencija: nazivi putanja] ... \end{konvencija}
\newenvironment{konvencija}[1][Konvencija]
  {\begin{exampleblock}{#1}}
  {\end{exampleblock}}

% ---------- Helper macros ----------
% Stupac tablice zadane širine, poravnat ulijevo (bez rastezanja i dijeljenja riječi).
\newcolumntype{L}[1]{>{\raggedright\arraybackslash}p{#1}}
\newcommand{\codefile}[1]{\texttt{#1}}
\newcommand{\kbd}[1]{\texttt{#1}}

% ---------- Code frame environment ----------
% Automatski dodaje fragile i shrink=5, koristi se umjesto \begin{frame}
% Primjer: \begin{codeframe}{Naslov} ... \end{codeframe}
% Primjer s opcijama: \begin{codeframe}[shrink=15]{Naslov} ... \end{codeframe}
\newenvironment{codeframe}[2][]
  {\begin{frame}[fragile,shrink=5,#1]{#2}}
  {\end{frame}}


\title[Predavanje 4: Autorizacija]{Programiranje za web \\ Predavanje 4: Autorizacija i vlasništvo nad zapisom}
\date{29.10.2026.}

\begin{document}

\frame{\titlepage}

% ------------------------------------------------------------
\begin{frame}{Plan predavanja}
\begin{enumerate}
  \item Što autorizacija odlučuje
  \item Provjera uloge
  \item Slojevi aplikacije
  \item Vlasništvo nad zapisom
  \item Odgovori i njihove posljedice
  \item Testiranje prava pristupa
\end{enumerate}

\vspace{0.6em}
Na kraju predavanja svaka ruta zna tko smije poslati zahtjev i nad čijim podatcima.
\end{frame}

% ------------------------------------------------------------
\begin{frame}{Ishodi učenja}
Nakon ovog predavanja student može:

\begin{itemize}
  \item razgraničiti provjeru uloge od provjere vlasništva i reći koju vrstu zaobilaženja svaka sprječava
  \item napisati ovisnost koja propušta samo određene uloge
  \item imenovati odgovornost svakoga sloja i reći koji od njih može odlučiti o pravu pristupa
  \item provesti provjeru vlasništva u sloju poslovnih pravila i obrazložiti zašto ne stoji u ruti
  \item razlikovati odbijanje zahtjeva od sužavanja odgovora
  \item obrazložiti izbor između odgovora \texttt{403} i \texttt{404}
  \item napisati test koji zahtjev nad tuđim zapisom odbija i prepoznati što padanje testa govori
\end{itemize}
\end{frame}

% ------------------------------------------------------------
\begin{frame}{Domena primjera}
Sustav prijava klubova na powerlifting natjecanje. Račun ima ulogu \texttt{admin} ili \texttt{club}. Račun kluba nosi oznaku svojega kluba, a organizator ne nosi nijednu.

\vspace{0.4em}
Iz prethodnog predavanja svaki zahtjev nosi token, a ovisnost \texttt{get\_current\_user} iz njega dobiva zapis o korisniku iz baze. Danas se odlučuje što taj korisnik smije.
\end{frame}

% ------------------------------------------------------------
\begin{frame}[fragile]{Dva zahtjeva koja izgledaju ispravno (I)}
Klub 12 prijavljen je u aplikaciju. Iz razvojne konzole preglednika šalje dva zahtjeva. Prvim pokušava stvoriti novi klub:

\begin{lstlisting}
POST /clubs
Authorization: Bearer <token kluba 12>

{"name": "Novi klub", "city": "Zadar", ...}
\end{lstlisting}

\vspace{0.4em}
Drugim traži podatke kluba 13:

\begin{lstlisting}
GET /clubs/13
Authorization: Bearer <token kluba 12>
\end{lstlisting}

\vspace{0.4em}
Token je valjan i nije istekao, obje adrese postoje, a tijelo je ispravnog oblika.
\end{frame}

% ------------------------------------------------------------
\begin{frame}{Dva zahtjeva koja izgledaju ispravno (II)}
Ovi zahtjevi ne smiju proći jer:

\vspace{0.9em}
\begin{itemize}
  \itemsep0.6em
  \item Prvi traži radnju koju smije samo organizator. Da bismo ga odbili, dovoljno je znati \textbf{tko} ga šalje.
  \item Drugi traži radnju koju klub smije, ali nad tuđim zapisom. Da bismo ga odbili, moramo pročitati \textbf{zapis} i usporediti ga s pošiljateljem.
\end{itemize}

\vspace{1.1em}
Prvo pitanje zovemo uloga, drugo vlasništvo. Ostatak predavanja odgovara na to gdje u aplikaciji stoji odluka o svakome od njih.
\end{frame}

% ============================================================
\sekcija{Što autorizacija odlučuje}{Granica između identiteta i prava, uloga i vlasništvo, modeli autorizacije.}

% ------------------------------------------------------------
\begin{frame}{Od identiteta do prava}
Znati tko šalje zahtjev nije isto što i znati smije li ga poslati. Primjer iz domene:

\vspace{0.5em}
\begin{itemize}
  \item Klub pokušava stvoriti drugi klub. Zahtjev je valjan i pošiljatelj je poznat, ali stvaranje klubova nije njegov posao nego organizatorov.
  \item Klub A dohvaća podatke kluba B. Za tu radnju ima pravo, ali samo nad zapisima svojega kluba.
\end{itemize}

\vspace{0.5em}
Prvi slučaj ovisi samo o tome tko je pošiljatelj. Drugi ovisi i o tome na koji se zapis zahtjev odnosi, pa se ne može odlučiti dok se ne zna o kojem je zapisu riječ.

\vspace{0.3em}
Zato autorizacija ima dva dijela, i oni se provode na različitim mjestima.
\end{frame}

% ------------------------------------------------------------
\begin{frame}{Uloga i vlasništvo}
\begin{center}
{\small
\setlength{\tabcolsep}{8pt}
\renewcommand{\arraystretch}{1.4}
\begin{tabular}{@{}lll@{}}
  \toprule
  \textbf{Provjera} & \textbf{Odgovara na} & \textbf{Sprječava} \\
  \midrule
  uloga      & smije li ovakvu radnju  & radnju iznad svojih prava \\
  vlasništvo & smije li nad ovim zapisom & dohvat tuđih podataka \\
  \bottomrule
\end{tabular}
}
\end{center}

\vspace{0.6em}
\textbf{Uloga} je svojstvo korisnika i ne ovisi o zahtjevu: klub nikada ne stvara klubove.

\vspace{0.3em}
\textbf{Vlasništvo} je odnos između korisnika i zapisa: isti korisnik smije mijenjati jedan zapis, a drugi ne, premda je radnja ista.

\vspace{0.4em}
Iz te razlike slijedi i mjesto provjere: uloga se može provjeriti prije nego bude dohvaćen zapis na koji se zahtjev odnosi, vlasništvo tek nakon toga.
\end{frame}

% ------------------------------------------------------------
\begin{frame}{Popis prava uz zapis (ACL)}
Uz svaki zapis stoji popis onih koji na njega imaju pravo i što smiju (engl. \emph{access control list}).

\vspace{1.1em}
Primjer je datoteka u operacijskom sustavu: vlasnik čita i piše, skupina samo čita, ostali ništa. Isto radi i dijeljeni dokument, gdje se pravo dodjeljuje po dokumentu, osobi po osobi.

\pause
\vspace{1.1em}
\begin{itemize}
  \itemsep0.6em
  \item \textbf{Prednost.} Pravo se dodjeljuje pojedinačno, bez izmjene koda.
  \item \textbf{Nedostatak.} Popis raste sa svakim zapisom i svakim korisnikom, a na pitanje što jedan korisnik sve smije odgovara se tek pretragom svih popisa.
\end{itemize}
\end{frame}

% ------------------------------------------------------------
\begin{frame}{Prava vezana uz ulogu (RBAC)}
Korisnik ima ulogu, a pravo je vezano uz ulogu, ne uz osobu (engl. \emph{role-based access control}).

\vspace{1.1em}
Primjer je uredništvo portala: uloga urednika smije objaviti članak, uloga novinara smije ga samo napisati i predati. Pravilo stoji na jednom mjestu i vrijedi za svakoga tko tu ulogu ima.

\pause
\vspace{1.1em}
\begin{itemize}
  \itemsep0.6em
  \item \textbf{Prednost.} Malo pravila pokriva mnogo korisnika, a uz rutu se vidi koja ih uloga smije pozvati.
  \item \textbf{Nedostatak.} Pravilo je isto za sve u ulozi, pa iznimka za jednoga traži novu ulogu. Vlasništvo uloga ne izražava: dva novinara imaju istu ulogu, a tuđi nacrt ne bi smjeli uređivati.
\end{itemize}
\end{frame}

% ------------------------------------------------------------
\begin{frame}{Prava izračunata iz svojstava (ABAC)}
Odluka se računa iz svojstava korisnika, zapisa i okolnosti u kojima zahtjev stiže (engl. \emph{attribute-based access control}).

\vspace{1.1em}
Primjer je pristup dokumentu u ustanovi: čitanje prolazi samo ako je razina provjere zaposlenika dovoljna za oznaku povjerljivosti dokumenta i ako zahtjev dolazi iz uredske mreže. Isti zaposlenik i isti dokument, a ishod ovisi o tome odakle pita.

\pause
\vspace{1.1em}
\begin{itemize}
  \itemsep0.6em
  \item \textbf{Prednost.} Obuhvaća pravila koja uloga ne može izraziti, a drže li se izvan koda, mijenjaju se bez nove inačice aplikacije.
  \item \textbf{Nedostatak.} Teže je odgovoriti tko što smije, jer odgovor ovisi o trenutku u kojem se pita.
\end{itemize}
\end{frame}

% ------------------------------------------------------------
\begin{frame}{Usporedba modela}
\begin{center}
{\small
\setlength{\tabcolsep}{8pt}
\renewcommand{\arraystretch}{1.4}
\begin{tabular}{@{}lll@{}}
  \toprule
  \textbf{Model} & \textbf{Odluka se donosi prema} & \textbf{Primjer} \\
  \midrule
  ACL  & popisu uz svaki zapis        & prava nad datotekom \\
  RBAC & ulozi korisnika              & urednik i novinar \\
  ABAC & svojstvima korisnika, zapisa i okolnosti & pravo koje ovisi o mreži \\
  \bottomrule
\end{tabular}
}
\end{center}

\vspace{0.9em}
U ovom se projektu koristi RBAC uz provjeru vlasništva: uloge su dvije, a pravila nije moguće promijeniti bez izmjene koda.
\end{frame}

% ============================================================
\sekcija{Provjera uloge}{Odakle dolazi uloga, kako se provjerava, uporaba u ruti.}

% ------------------------------------------------------------
\begin{frame}{Odakle dolazi uloga}
Uloga mora biti poznata pri svakom zahtjevu, a doći može s dvaju mjesta:

\vspace{0.4em}
\begin{center}
{\small
\setlength{\tabcolsep}{8pt}
\renewcommand{\arraystretch}{1.4}
\begin{tabular}{@{}lll@{}}
  \toprule
  \textbf{Izvor} & \textbf{Cijena} & \textbf{Dobitak} \\
  \midrule
  iz tokena & promjena kasni do isteka tokena & nijedan upit \\
  iz baze   & upit pri svakom zahtjevu        & promjena vrijedi odmah \\
  \bottomrule
\end{tabular}
}
\end{center}

\vspace{0.9em}
Oba su izvora u redu dok se uloga ne mijenja. Razlika se vidi tek kad se pravo oduzme.
\end{frame}

% ------------------------------------------------------------
\begin{frame}{Token nosi staro stanje}
Token nosi ono što je vrijedilo u trenutku izdavanja. Oduzme li se korisniku pravo, token izdan prije toga i dalje nosi staru ulogu dok ne istekne. Zapis u bazi nosi novu vrijednost odmah, ali za njega treba upit.

\vspace{0.8em}
\begin{konvencija}[Izbor, ne pravilo]
U našem se projektu uloga čita iz baze, jer \texttt{get\_current\_user} zapis o korisniku ionako dohvaća da provjeri je li račun aktivan; upit je time već plaćen. Aplikacije s vrlo velikim brojem zahtjeva odlučuju iz tokena i prihvaćaju da promjena prava kasni do njegova isteka.
\end{konvencija}
\end{frame}

% ------------------------------------------------------------
\begin{frame}[fragile]{Kako provjeravamo ulogu}
\begin{lstlisting}[language=python]
# app/core/deps.py
def require_role(*allowed_roles: str):
    def checker(current_user: User = Depends(get_current_user)) -> User:
        if current_user.role not in allowed_roles:
            raise AppError("forbidden", "Nemate dozvolu za ovu akciju", 403)
        return current_user
    return checker
\end{lstlisting}

\vspace{0.3em}
\texttt{require\_role} nije ovisnost nego funkcija koja je stvara. Poziv \texttt{require\_role("admin")} vraća funkciju \texttt{checker}, a nju ruta navodi u \texttt{Depends}. Tako jedan zapis pokriva sve kombinacije dopuštenih uloga.

\vspace{0.3em}
Unutar sebe \texttt{checker} navodi \texttt{get\_current\_user} kao svoju ovisnost, pa se provjera identiteta izvede prije provjere uloge. Rutama je time dovoljno navesti samo \texttt{require\_role}.
\end{frame}

% ------------------------------------------------------------
\begin{frame}[fragile]{Uporaba u ruti}
\begin{lstlisting}[language=python]
# app/routers/clubs.py
@router.get("/", response_model=list[ClubResponse])
async def list_clubs(
    db: AsyncSession = Depends(get_db),
    user: User = Depends(require_role("admin", "club")),
):
    ...

@router.post("/", response_model=ClubResponse, status_code=201)
async def create_club(
    body: ClubCreate,
    db: AsyncSession = Depends(get_db),
    admin: User = Depends(require_role("admin")),
):
    ...
\end{lstlisting}

\vspace{0.2em}
Tko rutu smije pozvati čita se iz njezina potpisa.
\end{frame}

% ------------------------------------------------------------
\begin{frame}[fragile]{Vježba: odbijena uloga u Swaggeru}
Najbrži način da se vidi odbijen zahtjev jest poslati ga.

\vspace{1.1em}
\begin{enumerate}
  \itemsep0.6em
  \item Prijavite se kao račun nekoga kluba, dakle s ulogom \texttt{club}:
\end{enumerate}

\begin{lstlisting}
POST /auth/login
{"username": "...", "password": "..."}
\end{lstlisting}

\begin{enumerate}
  \setcounter{enumi}{1}
  \itemsep0.6em
  \item Iz odgovora uzmite \texttt{access\_token} i unesite ga u \texttt{/docs}, gumbom \emph{Authorize}.
  \item Pošaljite \texttt{POST /clubs} s ispravnim tijelom i pogledajte odgovor.
\end{enumerate}

\vspace{0.8em}
Tijelo je ispravno, token je valjan, a zahtjev ne prolazi. Koji je kod odgovora i koja ga je točno ovisnost odbila?
\end{frame}

% ============================================================
\sekcija{Slojevi aplikacije}{Četiri sloja i njihove odgovornosti, zašto posao ne stoji u ruti, kako slojevi surađuju.}

% ------------------------------------------------------------
\begin{frame}{Koraci koje ponavlja svaka ruta}
Svaka ruta koja dira iste podatke prolazi kroz iste korake:

\vspace{0.9em}
\begin{itemize}
  \itemsep0.6em
  \item utvrditi tko šalje zahtjev
  \item provjeriti oblik poslanih podataka
  \item primijeniti pravilo
  \item pročitati ili upisati zapis
  \item sastaviti odgovor
\end{itemize}

\vspace{1.1em}
Kad bi svaka ruta sve to radila sama, isti bi kod stajao na više mjesta, a svaka bi se izmjena morala napraviti na svakom od njih.
\end{frame}

% ------------------------------------------------------------
\begin{frame}{Podjela posla po slojevima}
Posao se zato dijeli na slojeve. Svaki sloj obavlja jednu vrstu posla za sve rute koje ga trebaju i stoji u svojoj mapi:

\vspace{0.9em}
\begin{itemize}
  \itemsep0.6em
  \item \codefile{routers}: prima zahtjev i vraća odgovor
  \item \codefile{schemas}: opisuje oblik i tipove poslanih podataka
  \item \codefile{services}: provodi prava pristupa i poslovna pravila
  \item \codefile{repositories}: čita iz baze i upisuje u nju
\end{itemize}
\end{frame}

% ------------------------------------------------------------
\begin{frame}[fragile]{Sloj ruta: \codefile{routers}}
\begin{lstlisting}[language=python]
# app/routers/clubs.py
@router.get("/{club_id}", response_model=ClubResponse)
async def get_club(
    club_id: int,
    db: AsyncSession = Depends(get_db),
    user: User = Depends(require_role("admin", "club")),
):
    return await club_service.get_club(db, club_id, user)
\end{lstlisting}

\vspace{0.4em}
Dobro napisana ruta je kratka. Iz putanje uzima oznaku kluba, iz ovisnosti sesiju nad bazom i pošiljatelja, pa sve to predaje sloju \codefile{services}.

\vspace{1.1em}
Pravilo ne provodi sama, zapis ne čita sama i upit ne gradi sama. Njezin je posao prevesti zahtjev u poziv funkcije i odgovor te funkcije u HTTP odgovor.
\end{frame}

% ------------------------------------------------------------
\begin{frame}[fragile]{Isti posao u tri rute}
Dohvat kluba s pogreškom \texttt{404} potreban je trima rutama: čitanju, izmjeni i postavljanju nove lozinke. Bez sloja ispod sebe svaka bi ga radila sama:

\begin{lstlisting}[language=python]
# app/routers/clubs.py
async def get_club(club_id: int, ...):
    club = await club_repo.get_by_id(db, club_id)
    if not club:  # primjerak 1
        raise HTTPException(404)
async def update_club(club_id: int, ...):
    club = await club_repo.get_by_id(db, club_id)
    if not club:  # primjerak 2
        raise HTTPException(404)
async def reset_password(club_id: int, ...):
    club = await club_repo.get_by_id(db, club_id)
    if not club:  # primjerak 3
        raise HTTPException(404)
\end{lstlisting}
\end{frame}

% ------------------------------------------------------------
\begin{frame}{Sloj pravila: \codefile{services}}
Ispod ruta stoji sloj koji taj posao preuzima. Ovdje je odluka smije li se tražena radnja izvesti i što ona znači nad podatcima; u ovom projektu svakoj ruti odgovara jedna takva funkcija.

\vspace{0.9em}
\begin{itemize}
  \itemsep0.5em
  \item provodi poslovna pravila, među njima i provjeru vlasništva
  \item spaja više zapisa u jednu radnju: stvaranje kluba upisuje i u \texttt{clubs} i u \texttt{users}
  \item odbija zahtjev iznimkom \texttt{AppError}, koja nosi kod, poruku i statusni kod
\end{itemize}

\vspace{0.9em}
Sam odgovor ne sastavlja. \texttt{AppError} hvata zajednički rukovatelj i pretvara je u JSON, pa poruka o pogrešci ima isti oblik iz svake rute.
\end{frame}

% ------------------------------------------------------------
\begin{frame}[fragile]{Jedan dohvat, tri pozivatelja}
\begin{lstlisting}[language=python]
# app/services/club_service.py
async def _get_or_404(db: AsyncSession, club_id: int) -> Club:
    club = await club_repo.get_by_id(db, club_id)
    if not club:
        raise AppError("not_found", <@"Klub nije pronađen"@>, 404)
    return club

async def get_club(db, club_id, current_user):
    club = await _get_or_404(db, club_id)        # pozivatelj 1
async def update_club(db, club_id, body):
    club = await _get_or_404(db, club_id)        # pozivatelj 2
async def reset_password(db, club_id, new_password):
    await _get_or_404(db, club_id)               # pozivatelj 3
\end{lstlisting}

\vspace{0.2em}
Tri primjerka s prethodnog slajda postala su jedna funkcija i tri poziva. Poruka i kod pogreške isti su svuda, a izmjena se radi na jednom mjestu.
\end{frame}

% ------------------------------------------------------------
\begin{frame}[fragile]{Pravilo pomiješano s upitom}
Stvaranje kluba odbija se ako je naziv zauzet. Napiše li servis i sam upit, pravilo i dohvat stoje izmiješani:

\begin{lstlisting}[language=python]
result = await db.execute(select(Club).where(Club.name == body.name))
if result.scalar_one_or_none():
    raise AppError("duplicate", "Naziv kluba je zauzet", 409)
\end{lstlisting}

\vspace{0.8em}
S funkcijom koja dohvat obavlja umjesto njega, ostaje samo pravilo:

\begin{lstlisting}[language=python]
if await club_repo.get_by_name(db, body.name):
    raise AppError("duplicate", "Naziv kluba je zauzet", 409)
\end{lstlisting}

\vspace{0.8em}
Drugi se oblik čita kao rečenica: postoji li klub s tim nazivom, odbij zahtjev. Prvi kaže isto, ali tek nakon dvaju redaka o tome kako se do zapisa dolazi.
\end{frame}

% ------------------------------------------------------------
\begin{frame}[fragile]{Sloj pristupa bazi: \codefile{repositories}}
\begin{lstlisting}[language=python]
# app/repositories/club_repo.py
async def get_by_name(db: AsyncSession, name: str) -> Club | None:
    result = await db.execute(select(Club).where(Club.name == name))
    return result.scalar_one_or_none()

async def create(db: AsyncSession, club: Club) -> Club:
    db.add(club)
    <@\kom{flush, ne commit: transakciju potvrđuje get\_db kad zahtjev prođe}@>
    await db.flush()
    return club
\end{lstlisting}

\vspace{0.2em}
Sloj sadrži samo upite, ni jedno pravilo. Ne zna koja je uloga pošiljatelja ni čiji je zapis, pa ne može ni odlučiti smije li se izmjena izvesti.

\vspace{0.6em}
Zauzvrat se pravila čitaju bez poznavanja baze, a promjena pohrane ostaje na jednom mjestu.
\end{frame}

% ------------------------------------------------------------
\begin{frame}{Kako slojevi surađuju}
Primjer: organizator stvara novi klub.

\vspace{0.2em}
\begin{center}
\begin{tikzpicture}[
  node distance=7mm,
  box/.style={draw=ink, thick, rounded corners, fill=white, align=center,
              minimum height=9mm, minimum width=40mm, font=\footnotesize, text=textclr},
  vanj/.style={draw=lineclr, thick, dashed, rounded corners, fill=bgclr, align=center,
              minimum height=6.5mm, minimum width=18mm, font=\footnotesize, text=mutedclr},
  dolje/.style={-{Latex}, ink, thick},
  gore/.style={-{Latex}, ink2, thick},
  rub/.style={-{Latex}, ink2, thick, dashed},
  lbl/.style={font=\scriptsize\ttfamily, text=textclr, inner sep=2pt}
]
  \node[box] (r) {{\ttfamily routers}\\[1pt]{\scriptsize uloga \texttt{admin}, oblik tijela}};
  \node[box, below=of r] (s) {{\ttfamily services}\\[1pt]{\scriptsize naziv slobodan? hashiranje lozinke}};
  \node[box, below=of s] (p) {{\ttfamily repositories}\\[1pt]{\scriptsize dohvat i upis zapisa}};
  \node[vanj, below=of p] (d) {baza};
  \node[vanj, left=21mm of r] (q) {zahtjev};
  \node[vanj, right=21mm of r] (o) {odgovor};

  \draw[dolje] (q) -- (r) node[lbl, midway, above, yshift=2mm] {POST /clubs};
  \draw[gore] (r) -- (o) node[lbl, midway, above, yshift=2mm] {201 Created};

  \draw[dolje] ([xshift=-9mm]r.south) -- ([xshift=-9mm]s.north)
    node[lbl, midway, left, xshift=-1.5mm] {create\_club(db, body)};
  \draw[dolje] ([xshift=-9mm]s.south) -- ([xshift=-9mm]p.north)
    node[lbl, midway, left, xshift=-1.5mm] {create(db, Club(...))};
  \draw[rub] ([xshift=-9mm]p.south) -- ([xshift=-9mm]d.north)
    node[lbl, midway, left, xshift=-1.5mm] {INSERT INTO clubs};

  \draw[rub] ([xshift=9mm]d.north) -- ([xshift=9mm]p.south)
    node[lbl, midway, right, xshift=1.5mm] {dodijeljen id};
  \draw[gore] ([xshift=9mm]p.north) -- ([xshift=9mm]s.south)
    node[lbl, midway, right, xshift=1.5mm] {Club(id=7)};
  \draw[gore] ([xshift=9mm]s.north) -- ([xshift=9mm]r.south)
    node[lbl, midway, right, xshift=1.5mm] {Club(id=7)};
\end{tikzpicture}
\end{center}
\end{frame}

% ------------------------------------------------------------
\begin{frame}{Pravila suradnje}
\begin{itemize}
  \itemsep0.6em
  \item \textbf{Svaki sloj poziva samo onaj ispod sebe.} \codefile{services} ne zna da postoji \codefile{routers}, a \codefile{repositories} ne zna da postoje pravila.
  \item \textbf{Ruta pravilo navodi, ne provodi.} Uloga stoji u njezinu potpisu kao ovisnost, a izvodi je \texttt{require\_role}.
  \item \textbf{Sloj pravila ne gradi upite.} Poziva \texttt{club\_repo.get\_by\_name()} i ne zna kako je zapis pohranjen.
  \item \textbf{Sloj pristupa bazi ne zna tko šalje zahtjev.} Do njega ne dopiru ni uloga ni vlasništvo.
\end{itemize}

\vspace{0.9em}
Zato se svaki sloj može pozvati i bez onih iznad njega, \codefile{services} i bez pokrenutog poslužitelja.

\vspace{0.5em}
Cijena podjele je veći broj datoteka i jedan poziv više po zahtjevu. Isplati se čim isti posao pripada većem broju ruta.
\end{frame}

% ============================================================
\sekcija{Vlasništvo nad zapisom}{Zašto ne stoji u ruti, provjera u servisu, sužavanje popisa, zaobilaženje kroz oznaku u putanji.}

% ------------------------------------------------------------
\begin{frame}{Zašto vlasništvo ne stoji u ruti}
Provjera uloge gleda samo korisnika, pa se može izvesti prije nego zahtjev dođe do podataka. Provjera vlasništva to ne može, iz dvaju razloga:

\begin{itemize}
  \item \textbf{Zapis mora biti pročitan.} Ruta ga ne čita, pa ne može razlučiti zapis koji ne postoji od zapisa koji postoji, ali nije pošiljateljev.
  \item \textbf{Ruta vidi samo pošiljatelja.} Koja ga uloga smije pozvati stoji u njezinu potpisu. Čiji je zapis ne piše nigdje osim u zapisu samom.
\end{itemize}

\vspace{0.5em}
Vlasništvo je time poslovno pravilo, a ne svojstvo HTTP zahtjeva, pa stoji u sloju \codefile{services}, uz ostala pravila.
\end{frame}

% ------------------------------------------------------------
\begin{frame}[fragile]{Provjera u servisu}
\begin{lstlisting}[language=python]
# app/services/club_service.py
def ensure_club_access(current_user: User, club_id: int) -> None:
    if current_user.role == "club" and current_user.club_id != club_id:
        raise AppError(
            "forbidden", "Nemate pristup podatcima drugoga kluba", 403
        )

async def get_club(
    db: AsyncSession, club_id: int, current_user: User
) -> Club:
    club = await _get_or_404(db, club_id)
    ensure_club_access(current_user, club_id)
    return club
\end{lstlisting}

\vspace{0.4em}
\texttt{ensure\_club\_access} uspoređuje klub iz tokena s klubom iz putanje. Uvjet gleda samo ulogu \texttt{club}, jer organizator smije vidjeti svaki klub.
\end{frame}

% ------------------------------------------------------------
\begin{frame}[fragile]{Sužavanje popisa}
Dosadašnji zahtjevi imenovali su jedan zapis, pa se nad njim moglo pitati je li pošiljateljev. Zahtjev \texttt{GET /clubs} ne imenuje nijedan, pa nema što odbiti. Umjesto toga se bira što u odgovor ulazi:

\begin{lstlisting}[language=python]
# app/services/club_service.py
async def list_clubs(db: AsyncSession, current_user: User) -> list[Club]:
    if current_user.role == "admin":
        return await club_repo.get_all(db)
    club = await club_repo.get_by_id(db, current_user.club_id)
    return [club] if club else []
\end{lstlisting}

\vspace{0.3em}
Ista ruta, isti statusni kod, različit sadržaj. Klub ne dobiva \texttt{403} nego popis od jednog zapisa.

\vspace{0.3em}
To je i razlog zašto popis nikada ne smije dohvatiti sve pa filtrirati u odgovoru: \textbf{podatci koji ne pripadaju korisniku ne izlaze iz sloja koji ih čita}.
\end{frame}

% ------------------------------------------------------------
\begin{frame}[fragile]{Izmjena oznake u adresi}
Vratimo se drugom zahtjevu s početka predavanja:

\begin{lstlisting}
GET /clubs/13
Authorization: Bearer <token kluba 12>
\end{lstlisting}

\vspace{0.4em}
Sučelje tu adresu nikada ne otvori. Nastala je tako što je u postojećoj, \texttt{/clubs/12}, promijenjena jedna znamenka.

\vspace{1.1em}
Token je valjan, nije istekao i uloga je \texttt{club}, kakvu ruta i traži, pa provjera uloge prolazi. Bez usporedbe kluba iz tokena s klubom iz putanje, odgovor sadrži podatke kluba 13.

\vspace{1.1em}
Takav se propust naziva nesigurno izravno pozivanje na zapis (engl. \emph{insecure direct object reference}, IDOR).
\end{frame}

% ------------------------------------------------------------
\begin{frame}{Što ga ne sprječava}
\begin{itemize}
  \item \textbf{Sučelje koje tu adresu ne nudi.} Zahtjev ne mora doći iz sučelja. Alat naredbenog retka ili razvojna konzola preglednika šalju ga jednako dobro, s istim tokenom.
  \item \textbf{Oznaka koju je teško pogoditi.} Dovoljno ju je jednom vidjeti, a redni brojevi se ionako pogađaju redom. Nasumična oznaka usporava pretragu, ali ne mijenja ishod kad je oznaka poznata.
  \item \textbf{Provjera samo u putanji.} Oznaka zapisa ne mora stići samo adresom. Stigne li i u tijelu zahtjeva, provjera koja gleda samo putanju ostavlja tijelo nepokrivenim.
\end{itemize}

\vspace{0.5em}
Sprječava ga jedino usporedba pošiljatelja sa zapisom na koji se zahtjev odnosi, pri svakom zahtjevu i nad svakom oznakom koju je zahtjev donio.
\end{frame}

% ------------------------------------------------------------
\begin{frame}[fragile]{Vježba: redak koji sprječava IDOR}
Prijavite se kao jedan klub pa zatražite podatke drugoga:

\begin{lstlisting}
GET /clubs/<oznaka drugoga kluba>
Authorization: Bearer <token prvoga kluba>
\end{lstlisting}

\vspace{0.5em}
Odgovor je \texttt{403}. Odbio ga je jedan jedini redak koda. Nađite ga:

\vspace{0.6em}
\begin{enumerate}
  \itemsep0.5em
  \item Ubacite ispis prije i poslije poziva koje sumnjičite, ponovite zahtjev i u konzoli poslužitelja pogledajte dokle je izvođenje stiglo.
  \item Zakomentirajte redak koji ste našli i pošaljite isti zahtjev. Sada vraća \texttt{200} i podatke tuđega kluba.
  \item Vratite redak na mjesto i potvrdite da zahtjev opet dobiva \texttt{403}.
\end{enumerate}

\vspace{0.7em}
Ruta, token i uloga u sva su tri koraka bili isti.
\end{frame}

% ============================================================
\sekcija{Odgovori i njihove posljedice}{Razlika između \texttt{401} i \texttt{403}, izbor između \texttt{403} i \texttt{404}, provjere koje se ne oslanjaju jedna na drugu.}

% ------------------------------------------------------------
\begin{frame}{Kada \texttt{401}, a kada \texttt{403}}
\begin{center}
{\small
\setlength{\tabcolsep}{8pt}
\renewcommand{\arraystretch}{1.4}
\begin{tabular}{@{}lll@{}}
  \toprule
  \textbf{Kod} & \textbf{Stanje} & \textbf{Što klijentu preostaje} \\
  \midrule
  \texttt{401} & pošiljatelj nije utvrđen   & prijaviti se ili obnoviti token \\
  \texttt{403} & pošiljatelj je utvrđen     & ništa, zahtjev neće proći \\
  \bottomrule
\end{tabular}
}
\end{center}

\vspace{0.6em}
Razlika mijenja ponašanje klijenta. Na \texttt{401} klijent posegne za refresh tokenom i ponovi zahtjev. Vrati li poslužitelj \texttt{401} ondje gdje je riječ o nedostatku prava, klijent obnovi token, ponovi zahtjev, dobije isto, i korisnik ostane bez objašnjenja.

\vspace{0.3em}
U kodu je razlika vidljiva po tome tko podiže pogrešku: \texttt{get\_current\_user} vraća \texttt{401}, a \texttt{require\_role} i provjera vlasništva \texttt{403}.
\end{frame}

% ------------------------------------------------------------
\begin{frame}{Odati postojanje zapisa ili ne}
Klub traži podatke kluba koji postoji, ali nije njegov. Biramo odgovoriti s:

\begin{itemize}
  \item \texttt{403}: zapis postoji, ali nemaš pravo na njega
  \item \texttt{404}: kao da zapisa nema
\end{itemize}

\vspace{0.5em}
Prvi odgovor odaje da zapis postoji. Nizom zahtjeva s različitim oznakama dobiva se popis onoga što u sustavu postoji, bez ijednog dohvaćenog podatka.

\vspace{0.4em}
\begin{konvencija}[Izbor, ne pravilo]
Za podatke kojima samo postojanje nije tajna, \texttt{403} je jasniji jer korisniku kaže što se dogodilo. Ondje gdje je i postojanje zapisa osjetljivo, \texttt{404} odaje manje. U ovom projektu klubovi su javno poznati, pa se vraća \texttt{403}.
\end{konvencija}
\end{frame}

% ------------------------------------------------------------
\begin{frame}{Zaštita u više slojeva}
\begin{center}
{\small
\setlength{\tabcolsep}{8pt}
\renewcommand{\arraystretch}{1.35}
\begin{tabular}{@{}ll@{}}
  \toprule
  \textbf{Mjesto} & \textbf{Što zaustavlja} \\
  \midrule
  sučelje                & radnju koju korisnik ne treba vidjeti \\
  \codefile{routers}     & zahtjev bez tokena i pogrešnu ulogu \\
  \codefile{services}    & zahtjev nad tuđim zapisom \\
  baza                   & zapis koji krši ograničenje \\
  \bottomrule
\end{tabular}
}
\end{center}

\vspace{0.6em}
Ne smijemo se oslanjati na zaštitu preko sučelja: ono što ne prikaže i dalje se može poslati izravno na API.

\vspace{0.3em}
Provjera uloge ne zaustavlja dohvat tuđeg zapisa, a ograničenje u bazi ne zna ništa o tome tko je poslao zahtjev.
\end{frame}

% ------------------------------------------------------------
\begin{frame}{Zahtjevi koji ne bi smjeli proći}
Rutu obično isprobavamo kao korisnik koji na nju ima pravo, i tada sve radi. Izostane li provjera prava, ništa se ne mijenja: zahtjev i dalje prolazi, a odgovor izgleda jednako. Propust se pokaže tek kad isti zahtjev pošaljemo kao korisnik koji na njega nema pravo.

\vspace{0.5em}
Zato za svaku rutu moramo testirati:

\begin{itemize}
  \item bez tokena ili s neispravnim tokenom, uz očekivani \texttt{401}
  \item s tokenom korisnika pogrešne uloge, uz očekivani \texttt{403}
  \item s tokenom korisnika ispravne uloge, ali nad tuđim zapisom, uz očekivani \texttt{403}
\end{itemize}

\vspace{0.5em}
Treći je slučaj najlakše previdjeti, jer je zahtjev u svemu ispravan osim u jednom broju. Taj je propust ujedno i najčešći.
\end{frame}

% ============================================================
\sekcija{Testiranje prava pristupa}{Granica ručne provjere, oblik testa i priprema podataka, što znači kad test padne.}

% ------------------------------------------------------------
\begin{frame}{Granica ručne provjere}
Tri slučaja s prethodnog slajda vrijede za svaku rutu, a proći ih treba iznova nakon svake izmjene koda, ne samo kad pravilo pišemo.

\vspace{1.1em}
Provjera treba i onda kad pravilo ostaje isto. Izgubiti se može izmjenom koja s njim nema veze: poziv koji se premjesti, uvjet koji se proširi, funkcija koja dobije novoga pozivatelja. Nitko pravilo nije dirao, a više ne stoji.

\vspace{1.1em}
Nad klubovima je to pet ruta i tri računa: organizator, klub i drugi klub. Svaki prolaz traži dvije prijave i petnaestak zahtjeva. \textbf{Ručno to nije održivo.}
\end{frame}

% ------------------------------------------------------------
\begin{frame}{Krug jednoga testa}
Zato se zahtjev ne šalje rukom nego se zapisuje kao kod: program koji ga pošalje i provjeri odgovor. Takav se komad koda zove \textbf{test}.

\vspace{0.6em}
\begin{center}
\begin{tikzpicture}[
  node distance=6mm,
  box/.style={draw=ink, thick, rounded corners, fill=white, align=center,
              minimum height=8mm, minimum width=26mm, font=\footnotesize, text=textclr},
  str/.style={-{Latex}, ink, thick},
  povrat/.style={-{Latex}, ink2, thick, dashed},
  lbl/.style={font=\scriptsize, text=mutedclr, inner sep=2pt}
]
  \node[box] (a) {prazne tablice};
  \node[box, right=of a] (b) {pripremljeni zapisi};
  \node[box, right=of b] (c) {jedan zahtjev};
  \node[box, right=of c] (d) {tvrdnja o odgovoru};

  \draw[str] (a) -- (b);
  \draw[str] (b) -- (c);
  \draw[str] (c) -- (d);

  \draw[povrat] (d.south) -- ++(0,-7mm) -| (a.south);
  \node[lbl, below=9mm of b] {sljedeći test};
\end{tikzpicture}
\end{center}

\vspace{0.9em}
Svaki test prolazi isti krug. Prva dva koraka priprema \texttt{pytest}, alat koji testove pokreće; preostala dva čine sadržaj samoga testa.

\vspace{0.9em}
Krug je zatvoren: svaki test počinje od praznih tablica, pa ni jedan ne ovisi o onome što je ostavio prethodni. Redoslijed testova time ne utječe na ishod.
\end{frame}

% ------------------------------------------------------------
\begin{frame}[fragile]{Oblik testa}
\begin{lstlisting}[language=python]
# api/tests/test_clubs.py
async def test_club_cannot_get_other_club(
    client, club_and_user, club_b_and_user
):
    club_b, _ = club_b_and_user
    headers = await auth_header(client, "testclub", "klub123")
    <@\kom{IDOR: valjan token, ispravna uloga, tuđa oznaka u adresi}@>
    resp = await client.get(f"/clubs/{club_b.id}", headers=headers)
    assert resp.status_code == 403
    assert resp.json()["code"] == "forbidden"
\end{lstlisting}

\vspace{0.4em}
Test je jedan zahtjev i tvrdnja o odgovoru. Ovdje tvrdi dvoje: da je odgovor \texttt{403} i da u tijelu stoji \texttt{forbidden}, dakle da je zahtjev odbijen i zašto.
\end{frame}

% ------------------------------------------------------------
\begin{frame}[fragile]{Odakle dolaze podatci}
Test u zaglavlju navodi tri parametra, a nijedan mu nitko ne predaje. Ista ta imena nose funkcije u \texttt{conftest.py}:

\begin{lstlisting}[language=python]
# api/tests/conftest.py
@pytest.fixture
async def club_b_and_user(db: AsyncSession) -> tuple[Club, User]:
    ...

# api/tests/test_clubs.py
async def test_club_cannot_get_other_club(
    client, club_and_user, club_b_and_user
):
\end{lstlisting}

\vspace{0.4em}
Za svako ime pytest potraži funkciju označenu s \texttt{@pytest.fixture}, pozove je i rezultat preda testu. Takvu funkciju zovemo \emph{fixture}.

\vspace{0.8em}
\texttt{client} šalje zahtjeve aplikaciji bez pokrenutog poslužitelja.
\end{frame}

% ------------------------------------------------------------
\begin{frame}[fragile]{Fixture koji stvara klub}
\begin{lstlisting}[language=python]
# api/tests/conftest.py
@pytest.fixture
async def club_and_user(db: AsyncSession) -> tuple[Club, User]:
    club = Club(name="TestClub", city="Zadar")
    db.add(club)
    await db.flush()
    user = User(username="testclub", password_hash=fast_hash("klub123"),
                role="club", club_id=club.id)
    db.add(user)
    await db.commit()
    return club, user
\end{lstlisting}

\vspace{0.2em}
I sam fixture navodi parametar \texttt{db}, pa ga pytest popunjava na isti način.

\vspace{0.6em}
Test ovaj fixture navodi u zaglavlju, ali njegov rezultat ne upotrebljava. Treba mu samo ono što fixture ostavi u bazi: račun \texttt{testclub}, kojim se poslije prijavljuje pozivom \texttt{auth\_header}.
\end{frame}

% ------------------------------------------------------------
\begin{frame}[fragile]{Priprema u \codefile{conftest.py}}
Datoteku \texttt{conftest.py} pytest učitava sam. Fixturei iz nje dostupni su svim testovima u mapi, bez ijednog uvoza.

\begin{lstlisting}[language=python]
# api/tests/conftest.py
@pytest.fixture(scope="session", autouse=True)
def test_database():
    asyncio.run(_prepare_database())    # baza i tablice iz modela

@pytest.fixture(autouse=True)
async def clean_tables():
    tables = ", ".join(t.name for t in Base.metadata.sorted_tables)
    async with engine.begin() as conn:
        sql = f"TRUNCATE {tables} RESTART IDENTITY CASCADE"
        await conn.execute(text(sql))
\end{lstlisting}

\vspace{0.3em}
Nijedan ih test ne navodi u zaglavlju, a oba se izvedu.
\end{frame}

% ------------------------------------------------------------
\begin{frame}{Dvije oznake uz fixture}
\begin{itemize}
  \itemsep0.7em
  \item \texttt{autouse=True}: fixture se izvodi i kad ga nijedan test ne navede. Priprema koja vrijedi za sve ne mora stajati u svakom potpisu.
  \item \texttt{scope="session"}: izvodi se jednom za cijelo pokretanje. Bez te oznake fixture se izvodi prije svakoga testa, što je zadano ponašanje.
\end{itemize}

\vspace{1.1em}
Stvaranje baze i tablica dovoljno je jednom, pa \texttt{test\_database} nosi obje oznake. Pražnjenje tablica mora se dogoditi prije svakoga testa, pa \texttt{clean\_tables} nosi samo \texttt{autouse}.

\vspace{0.9em}
Redoslijed je time određen: jednom priprema baze, pa prazne tablice prije svakoga testa, pa fixturei koje test sam navede.
\end{frame}

% ------------------------------------------------------------
\begin{frame}[fragile]{Pokretanje testova}
Testovi rade nad bazom, pa se baza prvo podiže. Testna baza \texttt{pl\_reg\_test} odvojena je od razvojne i stvara se sama.

\begin{lstlisting}[language=bash]
docker compose up -d db     # u korijenu projekta
cd api
pytest                      # svi testovi
\end{lstlisting}

\vspace{0.6em}
Uz \texttt{-v} pytest imenuje svaki test i njegov ishod:

\begin{lstlisting}[language=bash]
pytest -v
tests/test_clubs.py::test_admin_creates_club_with_working_login PASSED
tests/test_clubs.py::test_create_club_without_token_is_401 PASSED
...
\end{lstlisting}
\end{frame}

% ------------------------------------------------------------
\begin{frame}[fragile]{Odabir testova}
Dok se piše jedno pravilo, cijeli skup nije potreban:

\begin{lstlisting}[language=bash]
pytest tests/test_clubs.py::test_club_cannot_get_other_club
pytest -k cannot            # izbor po imenu testa
pytest -x                   # stani na prvom padu
pytest -s                   # pusti ispis iz testova na zaslon
\end{lstlisting}

\vspace{0.9em}
Ono što test ispiše pytest inače zadržava i pokaže tek ako test padne, da prolaz ostane čitljiv. Opcija \texttt{-s} taj ispis pušta odmah, što pomaže dok se traži uzrok.

\vspace{0.9em}
Opcije se slažu: \texttt{pytest -vsx -k cannot} pokreće samo odabrane testove, imenuje ih, pušta njihov ispis i staje na prvom padu.
\end{frame}

% ------------------------------------------------------------
\begin{frame}{Vježba: što padanje testa kaže}
Testovi nad klubovima podijeljeni su u datoteci na dvije skupine, na dvije cjeline ovog predavanja: ulogu i vlasništvo. Ponovite zahvat s prethodne vježbe, ovaj put bez preglednika.

\vspace{0.9em}
\begin{enumerate}
  \itemsep0.5em
  \item Pokrenite \texttt{pytest} i provjerite da je sve zeleno.
  \item Zakomentirajte poziv \texttt{ensure\_club\_access} u \texttt{get\_club}. Prije nego pokrenete \texttt{pytest}, zapišite koliko testova očekujete da padne i kojih.
  \item Pokrenite \texttt{pytest -v} i usporedite s onim što ste zapisali. Zašto testovi uloge ostaju zeleni?
  \item Vratite poziv.
\end{enumerate}

\vspace{0.9em}
Na kraju: dodate li sutra rutu koja prima oznaku kluba, a zaboravite provjeru vlasništva, hoće li ijedan postojeći test pasti?
\end{frame}

% ------------------------------------------------------------
\begin{frame}{Sažetak}
\begin{itemize}
  \item Autorizacija odgovara na dva pitanja: smije li korisnik ovakvu radnju i smije li je nad ovim zapisom.
  \item Uloga se provjerava prije dohvata, u ovisnosti uz rutu; vlasništvo tek nakon što je zapis pročitan, u sloju pravila.
  \item Uloga koja ulazi u odluku dolazi iz baze, a ne iz tokena, pa promjena prava vrijedi odmah.
  \item Posao je podijeljen na slojeve, a o pravu pristupa odlučuje \codefile{services}, jedini sloj koji vidi i pošiljatelja i zapis.
  \item Zeleni testovi dokazuju da su pravila vrijedila za poslane zahtjeve, ne da stoje na svakoj ruti.
  \item Za popise nema što odbiti, pa se odgovor sužava: ista ruta, isti kod, različit sadržaj.
  \item \texttt{401} znači da se treba prijaviti, \texttt{403} da prijava neće pomoći; izbor između \texttt{403} i \texttt{404} odlučuje odaje li se postojanje zapisa.
\end{itemize}
\end{frame}

\end{document}
